An isomorphism between the two punctured affine lines, and its inverse, extend to morphisms $\mathbb{P}^1\to\mathbb{P}^1$ by Proposition 6.8. Their composites agree with the identity on dense open subsets, hence everywhere. Thus the extension is an automorphism, which bijects the omitted sets $\{P_1,\ldots,P_r,\infty\}$ and $\{Q_1,\ldots,Q_s,\infty\}$. Therefore $r=s$.

The converse fails. Choose $a\in k-\{0,1\}$ and choose $b$ outside $\{0,1,a,1-a,1/a,1/(1-a),a/(a-1),(a-1)/a\}$. The unordered quadruples $\{0,1,a,\infty\}$ and $\{0,1,b,\infty\}$ cannot be carried into one another by a fractional linear transformation: their cross-ratios differ even after all six possible changes from permuting the four points. By \exref{I.6.6}, $\mathbb{A}^1-\{0,1,a\}$ and $\mathbb{A}^1-\{0,1,b\}$ are therefore not isomorphic.
